\section{Dictators separate the subcritical regimes}
\label{sec:dictators}
The smallest nonconstant parity already shows that a lower bound
$d^{c k}$ cannot hold uniformly over all $k\ge1$ throughout
$\gamma<1$. Here and below the width tends to infinity at fixed
$d$, training sample, and exponent. We use the normalized Walsh
eigenvalues $\lambda_j$ of $K/2^d$.

\begin{lemma}[Exact ground states]\label{lem:dictator-ground}
The largest normalized eigenvalue is
$\lambda_*:=\lambda_0=\lambda_1$, and its eigenspace consists
exactly of affine functions on the cube. Moreover,
$\lambda_*\le1/(d+1)$. Among all vectors with $|z(x)|\ge a$
at every input, the minimum of $z^{\mathsf T}K^{-1}z/2$ is
$a^2/(2\lambda_*)$. Its minimizers are precisely
\[
 z\in\{\pm a,\ \pm a x_1,\ldots,\pm a x_d\}.
\]
\end{lemma}
\begin{proof}
The pairing and strict spectral ordering follow from
Lemma~\ref{lem:power}. For an elementary verification of strictness,
put $D=d+1$ and write $F(t)=\sum_{m\text{ odd}}c_mt^m$, where
every odd coefficient is positive. The degree-$j$ eigenvalue on the
full $D$-cube is $a_j=\sum_m c_mp_m(j)$, with $p_m(j)$ the
coordinate-flip walk probability of a specified distance-$j$ vertex.
Its exponential generating formula is
\[
 p_m(j)=\frac{m!}{D^m}[t^m](\sinh t)^j(\cosh t)^{D-j}.
\]
Subtracting the formula at $j+2$ leaves
$(\sinh t)^j(\cosh t)^{D-j-2}$, whose coefficients are
nonnegative. The coefficient at $m=j$ is positive. Thus
$a_j>a_{j+2}$ at every available odd pair, and restricting the
distinguished coordinate gives $\lambda_{2j}=\lambda_{2j+1}$.
Consequently the top eigenspace on the original cube has dimension
$d+1$ and consists of degrees zero and one. The normalized covariance
has trace $K(x,x)=1$, so $(d+1)\lambda_*\le1$.

For the uniform population norm,
\[
 z^{\mathsf T}K^{-1}z
 =\sum_U\frac{\widehat z(U)^2}{\lambda_{|U|}}
 \ge\frac{\|z\|_2^2}{\lambda_*}
 \ge\frac{a^2}{\lambda_*}.
\]
Equality requires both $|z(x)|=a$ everywhere and membership in the
top eigenspace. If $P=z/a=c_0+\sum_jc_jx_j$ takes values in
$\{-1,1\}$, then $P^2=1$. Its degree-two coefficients give
$c_ic_j=0$ for $i\ne j$, and its degree-one coefficients give
$c_0c_i=0$. Hence at most one of $c_0,c_1,\ldots,c_d$ is nonzero;
the constant coefficient of $P^2$ makes that coefficient $\pm1$.
The listed functions attain equality.
\end{proof}

\begin{theorem}[Logarithmic samples for a dictator]
\label{thm:dictator-fast}
Fix $L\ge2$, $a>0$, $0<\gamma<1$, and $0<\delta<1$.
For target $f(x)=x_1$, let $S$ consist of independent uniform
training inputs. If
\[
 n\ge\left\lceil\log_2\frac{2d}{\delta}\right\rceil,
\]
then, with probability at least $1-\delta$,
\[
 \lim_{N\to\infty}q_{N,S}(x)=1\qquad\text{for every }x\in\mathcal X.
\]
In particular the strict-$2/3$ confidence-failure fraction is zero
for all sufficiently large widths on this event.
\end{theorem}
\begin{proof}
The target is always a compatible ground-state sign. Its opposite
$-x_1$ is excluded by any training observation. Each of the other
$2d$ constants or signed coordinate functions agrees with the target
at a uniform input with probability $1/2$, and hence on all $n$
independent inputs with probability $2^{-n}$. A union bound shows
that the target is the unique compatible ground-state sign with
probability at least $1-2d\,2^{-n}\ge1-\delta$.

On this event, Lemma~\ref{lem:dictator-ground} makes the target the
unique compatible sign function attaining the global minimum
energy. Proposition~\ref{prop:rate} therefore puts posterior mass
tending to one on this sign function. This proves the assertion
simultaneously at all inputs, since the cube is finite at each width
limit. No formula for weights among tied minimizers is used.
\end{proof}

The contrast with the negative-exponent limit concerns actual
posterior probabilities as well. It does not identify confidence with
the sign of a Gaussian posterior mean.

\begin{theorem}[A universal linear-dimensional obstruction for
negative exponents]\label{thm:negative-linear}
Let $f:\mathcal X\to\{-1,1\}$ be any target and $S$ any training
sample of size $n$, allowing repetitions. Under the Gaussian prior
conditioned on the observed signs, set
\[
 F_\theta(S)=\mathbb P_x\{
       \mathbb P(f(x)G(x)>0\mid f(X_i)G(X_i)>0\ \forall i)
                  \le\theta\}.
\]
For every $1/2<\theta<1$,
\begin{equation}\label{eq:negative-linear}
 F_\theta(S)\ge
 1-\left(2+\frac1{\pi(\theta-1/2)^2}\right)\frac{n}{d+1}.
\end{equation}
Consequently, for every fixed $\gamma<0$, the actual neural-network
posterior with the stipulated global-margin convention satisfies
\begin{equation}\label{eq:negative-linear-width}
 \liminf_{N\to\infty}F_N(S)
 \ge1-\left(2+\frac{400}{9\pi}\right)\frac{n}{d+1}.
\end{equation}
In particular $n=o(d)$ leaves confidence failure $1-o(1)$, even
for a dictator or a constant target.
\end{theorem}
\begin{proof}
Remove repeated training inputs and let $s\le n$ be their number.
Write Gaussian regression on these values as
\[
 G(x)=R(x)+\varepsilon(x),\qquad
 R=K_{\mathcal X,S}K_{SS}^{-1}G_S,
 \qquad\sigma_x^2=\operatorname{Var}\varepsilon(x).
\]
The residual process is independent of $G_S$. Whiten the training
vector by $G_S=K_{SS}^{1/2}\xi$. The observed signs form a
homogeneous cone in the standard Gaussian vector $\xi$. Its radius
is independent of its direction, so conditioning on this cone
preserves $\mathbb E\|\xi\|^2=s$. With $\mathbb E_S^*$ denoting
expectation under that sign conditioning, the regression identity gives
\[
 \|R\|_K^2=\|\xi\|^2,\qquad
 \mathbb E_S^*\|R\|_2^2\le\lambda_*s\le\lambda_*n.
\]
These statements concern the second moment of the random regression
field, not just its posterior mean.

The prior covariance of $R$ is positive semidefinite, bounded by
$K$, and has rank at most $s$. Its normalized trace is at most
$s\lambda_*$. Since $\sigma_x^2=1-\operatorname{Var}R(x)$,
\[
 \mathbb P_x\{\sigma_x^2<1/2\}\le2n\lambda_*.
\]
For every remaining input, conditional Gaussian regression yields
\[
 q_S(x)=\mathbb E_S^*\Phi\!\left(\frac{f(x)R(x)}{\sigma_x}\right),
\]
where $\Phi$ is the standard normal distribution function. The
inequality $|\Phi(t)-1/2|\le |t|/\sqrt{2\pi}$ and
$\sigma_x\ge1/\sqrt2$ imply
\[
 |q_S(x)-1/2|
 \le\frac{\sqrt{\mathbb E_S^*R(x)^2}}{\sqrt\pi}.
\]
Thus $q_S(x)>\theta$ at such an input requires
$\mathbb E_S^*R(x)^2>\pi(\theta-1/2)^2$. Markov's inequality,
integrated over inputs, bounds the fraction of these inputs by
$n\lambda_*/[\pi(\theta-1/2)^2]$. Adding the variance exception
and using $\lambda_*\le1/(d+1)$ proves
\eqref{eq:negative-linear}.

Take $\theta=13/20<2/3$. For $\gamma<0$,
Proposition~\ref{prop:gp-limits} gives convergence of each of the
finitely many correct-label posterior probabilities to the
sign-conditioned Gaussian value. Every input with limiting value
at most $13/20$ therefore has finite-width value below $2/3$ for all
sufficiently large widths. This proves
\eqref{eq:negative-linear-width}. The case $n=0$ follows directly
from sign symmetry and also satisfies the displayed bounds.
\end{proof}

\section{An exact pairing of parity sizes}
\label{sec:projective-symmetry}
The following symmetry holds before taking any width limit.
\begin{proposition}[Projective-cube symmetry]
\label{prop:projective-symmetry}
At every width and exponent for which the conditioning is defined,
the iid-sample learning laws of coordinate parities of sizes
$0$ and $1$, of sizes $2$ and $3$, and so on, coincide whenever
both sizes are available in ambient dimension $d$. The equality
includes the distribution of the actual confidence-failure fraction.
\end{proposition}
\begin{proof}
Fix a coordinate $i$ and define an involution of the cube by
\[
 (T_ix)_i=x_i,\qquad (T_ix)_j=x_ix_j\quad(j\ne i).
\]
If $P_i$ swaps augmented coordinates $0$ and $i$, then
$(1,T_ix)=x_iP_i(1,x)$. For fixed weights $W$, replace only the first
matrix by $W'^1=W^1P_i$. Oddness of every hidden activation and of
the readout gives, inductively,
\[
 z_W(T_ix)=x_i z_{W'}(x).
\]
The laws of $W$ and $W'$ agree because the first-layer columns are
exchangeable. For a target $g$, define $g'(T_ix)=x_ig(x)$.
Absolute margins are invariant under this transformation, and the
training-label event for $(g',T_iS)$ becomes exactly the event for
$(g,S)$. Thus the correct-label posterior probabilities and their
failure fractions agree after the change of variables. The involution
preserves the uniform input distribution.

For $g=\chi_U$, $|U|=k$, writing $u=T_ix$ gives
\[
 g'(u)=u_i^{k+1}\prod_{j\in U\setminus\{i\}}u_j.
\]
If $k$ is odd and $i\in U$, this is a $(k-1)$-coordinate parity.
If $k$ is even and $i\notin U$, it is a $(k+1)$-coordinate parity.
Ordinary coordinate permutations identify all supports of a fixed
size, proving the assertion.
\end{proof}

These results refute a $d^{c k}$ lower bound asserted uniformly
including $k=1$ across the entire subcritical regime. They do not
refute a lower bound restricted to $k\to\infty$, and the symmetry
reduces parity size only within one adjacent even--odd pair.
